The key findings of this study are:

\begin{enumerate}
  \item \textbf{UMAP + MiniBatchKMeans} is the single best pipeline overall, achieving the
    highest Silhouette (0.541) and Adjusted Rand (0.494) simultaneously.
    UMAP + Birch wins on Calinski--Harabasz (35\,899) and Davies--Bouldin (0.658).
  \item \textbf{Non-linear DR (UMAP / t-SNE) is essential.} All internal metrics jump an
    order of magnitude compared to linear / reconstruction methods, because their embeddings
    push same-class points into tight, well-separated groups.
  \item \textbf{CNN-SAE + GMM} (ARI $\approx$ 0.463) is the best representation-learning
    approach, demonstrating that convolutional features capture genuine class structure even
    without explicit supervision.
  \item \textbf{DBSCAN consistently underperforms.} With auto-tuned hyperparameters it
    collapses most data into one cluster or noise in high-dimensional spaces; it only partly
    recovers in the compact t-SNE / UMAP regime.
  \item \textbf{Training cost is dominated by DR.} CNN-SAE ($\approx$6 min on GPU) and
    t-SNE / UMAP ($\approx$1.6 min) are the most expensive; PCA is essentially free.
    Non-linear DR also speeds up downstream clustering by producing lower-dimensional embeddings.
  \item \textbf{Inherent class overlap} (shirt / pullover / coat) is the hard ceiling; no
    unsupervised method fully resolves it without label information.
\end{enumerate}

\medskip
\noindent\textbf{Recommended pipeline:} UMAP + MiniBatchKMeans ($k=10$), offering the best
balance of clustering quality across all four metrics and acceptable runtime.

\medskip
\noindent\textbf{Was there a combination that achieved the best results on all metrics simultaneously?}
No. UMAP + MiniBatchKMeans achieves the best Silhouette (0.541) and Adjusted Rand (0.494),
while UMAP + Birch achieves the best Calinski--Harabasz (35\,899) and Davies--Bouldin (0.658).
No single combination dominates on all four metrics, though all top configurations share the
same DR method (UMAP), confirming it as the critical factor in clustering quality.
